% ======================================================================
% sections/limitations_future.tex
%
% Limitations and Future Work for the full paper:
%   Triple Humanoid 24/7 Adverse Event Oncology Trial Response Team:
%   4-Site Rotation (v0.7.0)
%
% Five subsections inside the single \section{Limitations and Future
% Work} block.
% ======================================================================

\begin{sectionaccent}

\subsection{Simulation Versus Real Robot Limitations}
\label{sec:limits-sim}

The paper is transparent about the boundary between the simulator
harness and a live Unitree H2 EDU body. The FDA RTCT submitter at
\fp{paper/codegen/src/fda\_rtct\_submitter.py} uses ed25519 signing
in production and a placeholder signature in simulation
\cite{fda_rtct_2026}. The 10 Hz in robot motion loop at
\fp{paper/codegen/src/robot\_loop.cpp} runs in a simulator harness;
it does not drive a live H2 EDU body \cite{unitree_h2_2026}. Sensor
frames produced by \fp{paper/codegen/src/sensor\_to\_xyz.py} are
synthetic in the current run. No real Unitree H2 EDU body executed
any of the broadcast commands documented in this paper. The FDA RTCT
pilot intake API is called only in simulation. Patient identifiers
are synthetic and have the form \fp{PAT-NET-001-PNNN}. These bounds
are inherent to a codegen tree designed to run on a high end
conventional server without specialized humanoid hardware.

\subsection{Per Site Data Identity Limitation}
\label{sec:limits-parity}

The 4 site decision streams at
\fp{paper/execution/data/site\_runtime\_runs/} share an identical
byte length of 62{,}330 bytes each. The parity is acceptable as a
schema and timing check: the 1 Hz cadence stamps every site at the
same rate on the same record schema. It is not yet a production
behavior. The next iteration must feed each site a site specific
adverse event arrival stream so the record content diverges by
site. A live run will produce per site record counts that differ
naturally as the arrival distribution diverges.

\subsection{Excluded Inputs}
\label{sec:limits-excluded}

The Notes section of the instructions tree explicitly excludes hours
56 through 83 of the national 24/7 trial dataset under
\fp{physical-ai-oncology-trials/new-trial/national-24-7-trial/extra-hours/}
\cite{kawchak_2026_humanoid_instructions,
kawchak_2026_physical_ai_oncology_trials_repo}. Only hours 00
through 55 are read into the current paper scope. The excluded
hours are a future inclusion roadmap item: once the supporting
analysis is ready, the hours 56 through 83 dataset can extend the
monitoring window beyond the current 168 hour scope.

\subsection{Missing Repository Update Instruction Directory}
\label{sec:limits-missing}

The Repository Structure block of the instructions tree names a
\fp{10-repository-update-instructions/} subdirectory that is not yet
present on disk \cite{kawchak_2026_humanoid_instructions}. This gap
is acknowledged. The next commit on the instructions tree should
create the directory and populate it with the top README,
releases.md, and CHANGELOG update instructions, mirroring the
pattern used by the 9 numbered subdirectories already present.

\subsection{Future Work Roadmap}
\label{sec:limits-roadmap}

Table \ref{tbl:future-roadmap} consolidates the future work roadmap
into a 7 row checklist. Each row pairs a roadmap item with the
anchor file that defines the current bound. The roadmap is the
hand off list from the technology layer to the governance layer.

\begin{table}[H]
\captionsetup{type=table}
\caption{Future work roadmap. Each item pairs with a project source
file that defines the current bound. Order is approximate: items
toward the top are blocked on technology readiness; items toward
the bottom are blocked on governance readiness.}
\label{tbl:future-roadmap}
\small
\begin{tabular}{>{\raggedright\arraybackslash}p{6cm}
                >{\raggedright\arraybackslash}p{9cm}}
\toprule
\textbf{Roadmap Item} & \textbf{Anchor File} \\
\midrule
Move from simulator harness to a live Unitree H2 EDU body with the kinematics chain as the live target & \fp{paper/codegen/src/kinematics.yaml} \\
Replace the placeholder signature with the live ed25519 signing chain & \fp{paper/codegen/src/fda\_rtct\_submitter.py} \\
Extend the adverse event arrival stream so the 4 site decision logs diverge in content & \fp{paper/execution/data/site\_runtime\_runs/} \\
Upgrade per site compute to NVIDIA Jetson AGX Thor 2070 TOPS where the workload demands it & \fp{paper/codegen/README.md} runtime section \\
Add the missing \fp{instructions/10-repository-update-instructions/} subdirectory & \fp{paper/instructions/README.md} \\
Include hour 56 through hour 83 extra hours dataset once supporting analysis is ready & \fp{physical-ai-oncology-trials/new-trial/national-24-7-trial/extra-hours/} \\
Enroll a real PAT-NET-001 cohort under an IRB protocol and a signed sponsor agreement & \fp{paper/codegen/src/fda\_rtct\_submitter.py} \\
\bottomrule
\end{tabular}
\end{table}

\subsubsection*{Real life feasibility insight.} The largest gap
between this paper and a real Phase 2 trial response team is not
the technology layer. It is the IRB, the sponsor, and the FDA RTCT
pilot governance. The technology layer is feasible today on the
hardware named in the codegen tree
\cite{unitree_h2_2026, nvidia_jetson_thor_2026,
anthropic_claude_opus_47_2026}. The governance layer is the work
that the paper invites the agencies, the sponsors, and the sites to
take on next.

\sectiontable{Limitations and Future Work planning matrix}{%
SimGap & L1 & 1 & robot loop, sub. & Sim vs real & 1 page & No live H2 run & Done \\
Parity & L2 & 2 & runtime runs & Identity check & 1 page & 62,330 bytes & Done \\
Excl. & L3 & 3 & instr Notes & Hours 56 to 83 & 0.5 page & By README & Done \\
Miss & L4 & 4 & instr struct. & 10-repo gap & 0.5 page & Not present & Done \\
Future & L5 & 5 & inputs, DR & 7 item table & 1 page & 7 items & Done \\
}

\end{sectionaccent}
